\section{第二章：矢量分析}

\begin{frame}{第二章：学习目标}
  \begin{itemize}
    \item 用分量、模与方向描述物理矢量。
    \item 掌握点积、叉积及其几何意义。
    \item 用 $\nabla$ 组织梯度、散度和旋度。
  \end{itemize}
  \begin{block}{核心思想}
    矢量把大小与方向同时编码；微分算子刻画场在空间中的局部变化。
  \end{block}
\end{frame}

\begin{frame}{矢量的分量表示}
  在三维直角坐标中，任意矢量可写为
  \[
    \mathbf{A}=A_x\mathbf{e}_x+A_y\mathbf{e}_y+A_z\mathbf{e}_z,
    \qquad
    |\mathbf{A}|=\sqrt{A_x^2+A_y^2+A_z^2}.
  \]
  \begin{columns}[T]
    \column{0.48\textwidth}
    \begin{block}{标量}
      只有大小，例如质量、温度、电荷。
    \end{block}
    \column{0.48\textwidth}
    \begin{block}{矢量}
      有大小和方向，例如位移、速度、力与电场。
    \end{block}
  \end{columns}
  \[
    \alpha\mathbf{A}=(\alpha A_x)\mathbf{e}_x+(\alpha A_y)\mathbf{e}_y+(\alpha A_z)\mathbf{e}_z.
  \]
  负的比例因子会反转矢量方向。
\end{frame}

\begin{frame}{加减法与点积}
  \begin{block}{分量运算}
    \[
      \mathbf{A}\pm\mathbf{B}=(A_x\pm B_x,\ A_y\pm B_y,\ A_z\pm B_z).
    \]
    加法对应平行四边形法则；$\mathbf{B}-\mathbf{A}$ 是从 $\mathbf{A}$ 的端点指向 $\mathbf{B}$ 的端点的矢量。
  \end{block}
  \begin{block}{点积}
    \[
      \mathbf{A}\cdot\mathbf{B}=|\mathbf{A}|\,|\mathbf{B}|\cos\theta
      =A_xB_x+A_yB_y+A_zB_z.
    \]
  \end{block}
  \begin{itemize}
    \item 点积为零 $\Longleftrightarrow$ 两个非零矢量正交。
    \item 在力学中，$W=\mathbf{F}\cdot\Delta\mathbf{r}$ 表示力做的功：只有沿位移方向的分量贡献功。
  \end{itemize}
\end{frame}

\begin{frame}{叉积与三重积}
  \[
    |\mathbf{A}\times\mathbf{B}|=|\mathbf{A}|\,|\mathbf{B}|\sin\theta.
  \]
  $\mathbf{A}\times\mathbf{B}$ 垂直于 $\mathbf{A}$ 和 $\mathbf{B}$，方向由右手定则确定；其模等于两矢量张成平行四边形的面积。
  \begin{block}{行列式写法}
    \[
      \mathbf{A}\times\mathbf{B}=
      \begin{vmatrix}
        \mathbf{e}_x&\mathbf{e}_y&\mathbf{e}_z\\
        A_x&A_y&A_z\\
        B_x&B_y&B_z
      \end{vmatrix}.
    \]
  \end{block}
  \[
    \mathbf{A}\cdot(\mathbf{B}\times\mathbf{C})
  \]
  称为标量三重积，其绝对值是平行六面体体积，正负号记录三矢量的取向。
\end{frame}

\begin{frame}{从标量场到矢量场}
  \begin{columns}[T]
    \column{0.48\textwidth}
    \begin{block}{梯度}
      对标量场 $\phi(x,y,z)$，
      \[
        \nabla\phi=\left(\frac{\partial\phi}{\partial x},
        \frac{\partial\phi}{\partial y},
        \frac{\partial\phi}{\partial z}\right).
      \]
      指向增长最快方向。
    \end{block}
    \column{0.48\textwidth}
    \begin{block}{散度与旋度}
      对矢量场 $\mathbf{F}$，
      \[
        \nabla\cdot\mathbf{F},\qquad \nabla\times\mathbf{F}.
      \]
      前者测量局部源汇，后者测量局部旋转。
    \end{block}
  \end{columns}
  \begin{alertblock}{重要恒等式}
    \[
      \nabla\times(\nabla\phi)=\mathbf{0},
      \qquad
      \nabla\cdot(\nabla\times\mathbf{F})=0.
    \]
  \end{alertblock}
\end{frame}

\begin{frame}{例：高斯定律与散度}
  电场 $\mathbf{E}$ 穿过闭曲面 $S$ 的通量为
  \[
    \Phi_E=\oiint_S\mathbf{E}\cdot\mathrm{d}\mathbf{S}.
  \]
  高斯定律给出
  \[
    \oiint_S\mathbf{E}\cdot\mathrm{d}\mathbf{S}=\frac{Q_{\mathrm{enc}}}{\varepsilon_0}.
  \]
  利用散度定理，
  \[
    \oiint_S\mathbf{E}\cdot\mathrm{d}\mathbf{S}
    =\iiint_V(\nabla\cdot\mathbf{E})\,\mathrm{d}V,
  \]
  因而得到微分形式
  \[
    \nabla\cdot\mathbf{E}=\frac{\rho}{\varepsilon_0}.
  \]
  \begin{block}{读法}
    电荷密度 $\rho$ 是电场发散的局部来源：正电荷是源，负电荷是汇。
  \end{block}
\end{frame}
